% ==========================================================================
% GENERATED by scripts/split_paper.py. Do not edit this file.
%
% Source: imported_papers/Neural_Computation__CNA/NECO-D-24-00304-Source.tex
% Part 15: Autoencoder Training
%
% Edits belong upstream, in the published paper. To change what the
% chapter includes, change imports/cna.toml and re-run
% the splitter.
% ==========================================================================
\section{Autoencoder Training}\label{sec:autoencoder_params}
We compare our CNA model to an ANN, specifically an autoencoder \citep{Hinton2006}, which is a widely used self-supervised technique for generating input representations. An autoencoder is particularly suited for comparison because it includes a built-in decoder, allowing for feature reconstruction and visualization, which is required for feature comparison. To ensure a fair comparison, we also employ convolutional filters \citep{lecun_backpropagation_1989} to the autoencoder and incorporate state-of-the-art design principles (we even permit more data and feature channels to the autoencoder).

The autoencoder is trained to reconstruct the feature activations, $\boldsymbol{y}^{(S1)}$, using the same training dataset as our CNA model. It consists of an 8-layer architecture, with an encoder composed of convolutional layers with a kernel size of $3$ and a stride of $2$, with progressively increasing channel counts of $32$, $64$, $128$, and $256$ channels. The decoder mirrors this structure, using transposed convolutional layers with $256$, $128$, $64$, and $32$ channels, respectively. ReLU activations are applied after each layer. Consequently, the autoencoder is a relatively large model with $777$k parameters, utilizing up to $256$ channels, whereas the CNA model operates with only $44$ channels.

The autoencoder is trained by minimizing the L2 norm between the input features and their reconstructed counterparts, using the Adam optimizer \citep{Kingma_Ba_2015} with a learning rate of $1 \times 10^{-4}$, beta parameters of $(0.9, 0.999)$, and a weight decay of $1 \times 10^{-8}$.

To provide a robust comparison, we increase the number of training samples for the autoencoder. Rather than training for 100 epochs with 300 samples per epoch as done for the CNA model, we train the autoencoder for $200$ epochs, using $10,000$ images per epoch. This results in a significantly larger training set but ensures that the autoencoder achieves a highly accurate reconstruction. The autoencoder is trained with a batch size of $512$.
